
\documentclass{article}
\usepackage{colortbl}
\usepackage{makecell}
\usepackage{multirow}
\usepackage{supertabular}

\begin{document}

\newcounter{utterance}

\centering \large Interaction Transcript for game `cladder', experiment `full\_v1.5\_default', episode 9125 with refineonly{-}Qwen3.5{-}4B.
\vspace{24pt}

{ \footnotesize  \setcounter{utterance}{1}
\setlength{\tabcolsep}{0pt}
\begin{supertabular}{c@{$\;$}|p{.15\linewidth}@{}p{.15\linewidth}p{.15\linewidth}p{.15\linewidth}p{.15\linewidth}p{.15\linewidth}}
    \# & \multicolumn{2}{c}{Player} && \multicolumn{2}{c}{Game Master} \\
    \hline

    \theutterance \stepcounter{utterance}  
    & & & \multicolumn{4}{p{0.6\linewidth}}{
        \cellcolor[rgb]{0.9,0.9,0.9}{
            \makecell[{{p{\linewidth}}}]{
                \texttt{\tiny{[P1$\langle$GM]}}
                \texttt{You are an expert in causal inference. The following question is not a typical commonsense query, but rather a meticulously designed question created by a professor specializing in causal inference, intended to assess the students' mastery of the course content.} \\
\texttt{Imagine a self{-}contained, hypothetical world with only the following conditions, and without any unmentioned factors or causal relationships: Kraz has a direct effect on pexu and rukz. Pexu has a direct effect on hwax. Hwax has a direct effect on rukz. Kraz is unobserved. For those who are not pexu and are not hwax, the probability of rukz is 19\%. For those who are not pexu and are hwax, the probability of rukz is 29\%. For those who are pexu and are not hwax, the probability of rukz is 20\%. For those who are pexu and are hwax, the probability of rukz is 31\%. For those who are not pexu, the probability of hwax is 16\%. For those who are pexu, the probability of hwax is 90\%. For those who are pexu, would it be less likely to see rukz if the individual was not pexu?} \\
\texttt{Start your answer with 'Yes' or 'No', followed by additional reasoning or evidence to support your explanation.} \\
            }
        }
    }
     \\ \\

    \theutterance \stepcounter{utterance}  
    & & \multicolumn{4}{p{0.6\linewidth}}{
        \cellcolor[rgb]{0.95,0.95,0.95}{
            \makecell[{{p{\linewidth}}}]{
                \texttt{\tiny{[P1$\leftrightarrow$P]}}
                \texttt{{[}DRAFT{]}} \\
\texttt{Yes} \\
\\ 
\texttt{**Reasoning:**} \\
\\ 
\texttt{To determine whether it would be "less likely to see rukz if the individual was not pexu," we must compare the marginal probability of \$rukz\$ given \$pexu\$ versus the marginal probability of \$rukz\$ given \$\textbackslash{}neg pexu\$.} \\
\\ 
\texttt{**Step 1: Calculate \$P(rukz | pexu)\$**} \\
\texttt{We are given the conditional probabilities of \$rukz\$ based on the joint states of \$pexu\$ and \$hwax\$. We also need the conditional probability of \$hwax\$ given \$pexu\$.} \\
\texttt{*   \$P(rukz | pexu, \textbackslash{}neg hwax) = 20\textbackslash{}\%\$} \\
\texttt{*   \$P(rukz | pexu, hwax) = 31\textbackslash{}\%\$} \\
\texttt{*   \$P(hwax | pexu) = 90\textbackslash{}\%\$ (which implies \$P(\textbackslash{}neg hwax | pexu) = 10\textbackslash{}\%\$)} \\
\\ 
\texttt{Using the law of total probability:} \\
\texttt{\$\$P(rukz | pexu) = P(rukz | pexu, \textbackslash{}neg hwax) \textbackslash{}cdot P(\textbackslash{}neg hwax | pexu) + P(rukz | pexu, hwax) \textbackslash{}cdot P(hwax | pexu)\$\$} \\
\texttt{\$\$P(rukz | pexu) = (0.20 \textbackslash{}times 0.10) + (0.31 \textbackslash{}times 0.90)\$\$} \\
\texttt{\$\$P(rukz | pexu) = 0.02 + 0.279 = 0.299 \textbackslash{}text\{ (or \} 29.9\textbackslash{}\%)\$\$} \\
\\ 
\texttt{**Step 2: Calculate \$P(rukz | \textbackslash{}neg pexu)\$**} \\
\texttt{We are given the probabilities for individuals who are not \$pexu\$.} \\
\texttt{*   \$P(rukz | \textbackslash{}neg pexu, \textbackslash{}neg hwax) = 19\textbackslash{}\%\$} \\
\texttt{*   \$P(rukz | \textbackslash{}neg pexu, hwax) = 29\textbackslash{}\%\$} \\
\texttt{*   \$P(hwax | \textbackslash{}neg pexu) = 16\textbackslash{}\%\$ (which implies \$P(\textbackslash{}neg hwax | \textbackslash{}neg pexu) = 84\textbackslash{}\%\$)} \\
\\ 
\texttt{Using the law of total probability:} \\
\texttt{\$\$P(rukz | \textbackslash{}neg pexu) = P(rukz | \textbackslash{}neg pexu, \textbackslash{}neg hwax) \textbackslash{}cdot P(\textbackslash{}neg hwax | \textbackslash{}neg pexu) + P(rukz | \textbackslash{}neg pexu, hwax) \textbackslash{}cdot P(hwax | \textbackslash{}neg pexu)\$\$} \\
\texttt{\$\$P(rukz | \textbackslash{}neg pexu) = (0.1} \\
            }
        }
    }
    & \\ \\

    \theutterance \stepcounter{utterance}  
    & \multicolumn{4}{p{0.6\linewidth}}{
        \cellcolor[rgb]{0.561,0.737,0.561}{
            \makecell[{{p{\linewidth}}}]{
                \texttt{\tiny{[P1$\rangle$GM]}}
                \texttt{Yes} \\
\\ 
\texttt{**Reasoning:**} \\
\\ 
\texttt{To determine whether it would be "less likely to see rukz if the individual was not pexu," we must compare the marginal probability of \$rukz\$ given \$pexu\$ versus the marginal probability of \$rukz\$ given \$\textbackslash{}neg pexu\$.} \\
\\ 
\texttt{**Step 1: Calculate \$P(rukz | pexu)\$**} \\
\texttt{We are given the conditional probabilities of \$rukz\$ based on the joint states of \$pexu\$ and \$hwax\$. We also need the conditional probability of \$hwax\$ given \$pexu\$.} \\
\texttt{*   \$P(rukz | pexu, \textbackslash{}neg hwax) = 20\textbackslash{}\%\$} \\
\texttt{*   \$P(rukz | pexu, hwax) = 31\textbackslash{}\%\$} \\
\texttt{*   \$P(hwax | pexu) = 90\textbackslash{}\%\$ (which implies \$P(\textbackslash{}neg hwax | pexu) = 10\textbackslash{}\%\$)} \\
\\ 
\texttt{Using the law of total probability:} \\
\texttt{\$\$P(rukz | pexu) = P(rukz | pexu, \textbackslash{}neg hwax) \textbackslash{}cdot P(\textbackslash{}neg hwax | pexu) + P(rukz | pexu, hwax) \textbackslash{}cdot P(hwax | pexu)\$\$} \\
\texttt{\$\$P(rukz | pexu) = (0.20 \textbackslash{}times 0.10) + (0.31 \textbackslash{}times 0.90)\$\$} \\
\texttt{\$\$P(rukz | pexu) = 0.02 + 0.279 = 0.299 \textbackslash{}text\{ (or \} 29.9\textbackslash{}\%)\$\$} \\
\\ 
\texttt{**Step 2: Calculate \$P(rukz | \textbackslash{}neg pexu)\$**} \\
\texttt{We are given the probabilities for individuals who are not \$pexu\$.} \\
\texttt{*   \$P(rukz | \textbackslash{}neg pexu, \textbackslash{}neg hwax) = 19\textbackslash{}\%\$} \\
\texttt{*   \$P(rukz | \textbackslash{}neg pexu, hwax) = 29\textbackslash{}\%\$} \\
\texttt{*   \$P(hwax | \textbackslash{}neg pexu) = 16\textbackslash{}\%\$ (which implies \$P(\textbackslash{}neg hwax | \textbackslash{}neg pexu) = 84\textbackslash{}\%\$)} \\
\\ 
\texttt{Using the law of total probability:} \\
\texttt{\$\$P(rukz | \textbackslash{}neg pexu) = P(rukz | \textbackslash{}neg pexu, \textbackslash{}neg hwax) \textbackslash{}cdot P(\textbackslash{}neg hwax | \textbackslash{}neg pexu) + P(rukz | \textbackslash{}neg pexu, hwax) \textbackslash{}cdot P(hwax | \textbackslash{}neg pexu)\$\$} \\
\texttt{\$\$P(rukz | \textbackslash{}neg pexu) = (0.1} \\
            }
        }
    }
    & & \\ \\

    \theutterance \stepcounter{utterance}  
    & & & \multicolumn{2}{p{0.3\linewidth}}{
        \cellcolor[rgb]{0.9,0.9,0.9}{
            \makecell[{{p{\linewidth}}}]{
                \texttt{\tiny{[GM$|$GM]}}
                \texttt{Yes} \\
            }
        }
    }
    & & \\ \\

    \theutterance \stepcounter{utterance}  
    & & & \multicolumn{2}{p{0.3\linewidth}}{
        \cellcolor[rgb]{0.9,0.9,0.9}{
            \makecell[{{p{\linewidth}}}]{
                \texttt{\tiny{[GM$|$GM]}}
                \texttt{yes} \\
            }
        }
    }
    & & \\ \\

    \theutterance \stepcounter{utterance}  
    & & & \multicolumn{2}{p{0.3\linewidth}}{
        \cellcolor[rgb]{0.9,0.9,0.9}{
            \makecell[{{p{\linewidth}}}]{
                \texttt{\tiny{[GM$|$GM]}}
                \texttt{game\_result = WIN} \\
            }
        }
    }
    & & \\ \\

\end{supertabular}
}

\end{document}
